\section{All Palindromes}
\label{sec:cses-all-palindromes}

\problemstatement{Given a string, count the total number of palindromic
substrings (each occurrence counted, for every start/end pair that forms a
palindrome).}

\begin{patternbox}
\emph{``Count all palindromic substrings''} is again \textbf{Manacher's
algorithm}: the radius at each centre directly gives how many palindromes
are centred there, so summing the radii counts every palindromic substring
in $O(n)$.
\end{patternbox}

\subsection*{Radii are counts}

For an odd centre at character $i$ with maximal radius $d_1[i]$ (counting
the centre itself), there are exactly $d_1[i]$ odd-length palindromes
centred at $i$: lengths $1,3,\ldots,2d_1[i]-1$. Likewise each even centre
between two characters contributes $d_2[i]$ even-length palindromes. So
the total number of palindromic substrings is
$\sum_i d_1[i]+\sum_i d_2[i]$. Manacher computes both radius arrays in
linear time, and the sum is the answer.

\begin{ideabox}[Every Nested Palindrome Counts Once]
A palindrome of radius $k$ around a centre contains, at the same centre,
palindromes of every smaller radius -- $k$ of them. Adding the radius at
each centre therefore counts every (centre, length) palindrome exactly
once, with no double counting across centres.
\end{ideabox}

\subsection*{Algorithm}

\begin{enumerate}
  \item Run Manacher on the sentinel-transformed string to get radii
        $P[i]$.
  \item In the transformed indexing, each centre contributes
        $\lceil P[i]/2\rceil$ palindromic substrings; sum over all
        centres (a single combined pass counts both odd and even).
  \item Output the total.
\end{enumerate}

\subsection*{Dry run}

\dryrun{$s=$``aaa''.}

\begin{itemize}
  \item Single characters: $3$ palindromes. Length-2 ``aa'': $2$.
        Length-3 ``aaa'': $1$.
  \item Total $=3+2+1=6$, which equals the sum of Manacher radii.
\end{itemize}

\subsection*{C++ implementation}

\begin{lstlisting}
#include <bits/stdc++.h>
using namespace std;

int main() {
    string s;
    cin >> s;
    string t = "^";
    for (char c : s) { t += '#'; t += c; }
    t += "#$";

    int n = t.size();
    vector<int> P(n, 0);
    int c = 0, r = 0;
    long long total = 0;
    for (int i = 1; i < n - 1; i++) {
        if (i < r) P[i] = min(r - i, P[2 * c - i]);
        while (t[i + P[i] + 1] == t[i - P[i] - 1]) P[i]++;
        if (i + P[i] > r) { c = i; r = i + P[i]; }
        total += (P[i] + 1) / 2;                  // palindromes at this center
    }
    cout << total << "\n";
    return 0;
}
\end{lstlisting}

\begin{complexitybox}
\textbf{Time Complexity.} $O(n)$ -- one Manacher pass. \textbf{Space
Complexity.} $O(n)$ for the transformed string and radii.
\end{complexitybox}

\begin{takeawaybox}
The radius Manacher computes at each centre \emph{is} the count of
palindromes there, so counting all palindromic substrings is a single sum.
For counting \emph{distinct} palindromes instead, an Eertree (palindromic
tree) is the specialised structure -- but for total occurrences, Manacher
suffices.
\end{takeawaybox}
